\section{A bounded-local reversible rewrite model}
\label{sec:model}

This section isolates exactly the hypotheses used later.  The model is broad
enough for the fixed-dart Reidemeister implementation and narrow enough that
commutation arguments are valid without hidden geometric assumptions.

\subsection{States, actions, and stable sites}

Fix a finite set $V$ of \emph{sites}, with $|V|\leq c_0N$ for an input-size
parameter $N$.  A state belongs to a set $\Omega_N$.  There are two kinds of
local operation.

\begin{definition}[Neutral action]
A neutral action instance $a$ legal at $x\in\Omega_N$ is a reversible transition
\[
  x \xrightarrow{a} a x
\]
with a footprint $\Phi_x(a)\subseteq V$.  We assume
$|\Phi_x(a)|\leq f$, where $f$ is independent of $N$.
\end{definition}

\begin{definition}[Terminal reduction]
A terminal reduction instance $r$ legal at $x$ is a local decreasing operation
with footprint $\Psi_x(r)\subseteq V$, with
$|\Psi_x(r)|\leq g$.  During an unlocking search we only ask whether $r$ is
legal; the reduction is applied after a trace succeeds.
\end{definition}

For RIII, the neutral action preserves the number of crossings.  RI/RII are the
terminal reductions.  Stable sites are fixed input darts, so a crossing can be
moved combinatorially without renaming its darts.

\subsection{Strong locality}

The main axiom is a two-sided commutation diamond.  It is stated for executed
action instances, because a later move need not be syntactically visible before
an earlier local update.

\begin{definition}[Strong disjoint commutation]
A rewrite system has strong disjoint commutation when the following hold.
\begin{enumerate}[label=(L\arabic*)]
\item \textbf{Backward and forward diamond.}  Suppose
      $x\xrightarrow{a}y\xrightarrow{b}z$, and the residual footprints of $a$
      and $b$ are disjoint.  Then $b$ is legal at $x$, $a$ is legal after $b$,
      both orders reach the same state, and the two footprints are unchanged by
      swapping the order.
\item \textbf{Terminal locality.}  If
      $x\xrightarrow{a}y$ and a terminal reduction $r$ is legal at one endpoint
      with footprint disjoint from $\Phi_x(a)$, then its corresponding residual
      is legal at the other endpoint with the same footprint.
\item \textbf{Reversibility.}  Every neutral action has a local inverse whose
      footprint is the same bounded set of sites.
\end{enumerate}
\end{definition}

The backward half of (L1) is essential.  It says that a move observed after an
unrelated action was already legal before that action.  This is what permits a
birth action to commute left.

\begin{remark}
For a concrete implementation, one can use a conservative footprint larger than
the minimal topological support.  This can only classify more actions as
interacting.  The search becomes less sharply pruned, but all commutation claims
remain safe.
\end{remark}

\subsection{Incidence bounds}

We separate the number of moves available globally from the number meeting a
small active region.

\begin{definition}[Bounded incidence]
A rule family has constants $A,B>0$ when, for every state $x$,
\begin{align*}
  |\{a:a\text{ neutral and legal at }x\}| &\leq AN, \\
  |\{a:a\text{ neutral and legal at }x,
      \ \Phi_x(a)\cap S\neq\varnothing\}| &\leq B|S|
\end{align*}
for every $S\subseteq V$.
\end{definition}

The first bound permits $O(N)$ possible locations for a new independent branch.
The second says that continuation from a region of $O(k)$ sites has only $O(k)$
choices.  In a finite local rule system on a bounded-degree combinatorial
structure, such constants are natural.

\subsection{Unlocking traces}

Let
\[
 x_0\xrightarrow{a_1}x_1\xrightarrow{a_2}\cdots
 \xrightarrow{a_k}x_k
\]
be a neutral trace.

\begin{definition}[First-unlocking trace]
The trace is first-unlocking when no terminal reduction is legal at
$x_0,\ldots,x_{k-1}$ and at least one terminal reduction $r$ is legal at $x_k$.
Its length is $k$.  It is \emph{shortest} when no first-unlocking trace from
$x_0$ has smaller length.
\end{definition}

A search can stop at the first terminal reduction.  Therefore any trace that
exposes a reduction earlier than expected is automatically replaced by its
shorter prefix.

\begin{definition}[Birth and birth number]
Write $\Phi_i=\Phi_{x_{i-1}}(a_i)$.  Action $a_i$ is a \emph{birth} when
\[
  \Phi_i\cap\bigcup_{j<i}\Phi_j=\varnothing.
\]
The birth number $\beta(\tau)$ is the number of birth actions in $\tau$.
\end{definition}

The first action is always a birth.  A birth is not necessarily topologically
remote forever: later actions may connect its footprint to other branches.

\subsection{The action-support graph}

Fix a terminal reduction $r$ legal at $x_k$.  The action-support graph
$H(\tau,r)$ has vertices $1,\ldots,k,\star$.  Vertex $i$ represents $a_i$ and
$\star$ represents $r$.  Put an edge between two action vertices when their
executed footprints intersect, and an edge $i\star$ when $\Phi_i$ intersects
the terminal footprint $\Psi_{x_k}(r)$ after the canonical residual
identification supplied by strong locality.

In the fixed-site situations used here, the residual identification is literal:
the same dart labels are compared throughout.  The graph records causal
noncommutation, not chronological adjacency.

\begin{example}[Two branches before a join]
Let actions $A$ and $B$ have disjoint footprints.  Let $C$ overlap both, and let
the terminal reduction overlap $C$.  The chronological word $ABC$ has two
support components after its second letter, but the full action-support graph is
connected.  This is the smallest obstruction to the naive claim that every
shortest trace can be given connected prefixes.
\end{example}
